\section*{Chapter \ref{ch:g6:pure-substances-and-mixtures} --- Pure Substances and Mixtures}

\begin{solution}{exo:g6:pure-substances-and-mixtures:1}
\omterm{def:g6:pure-substances-and-mixtures:pure-substance}{Pure substances}: distilled water, the block of pure iron, the sugar
crystal. \omterm{def:g6:pure-substances-and-mixtures:mixture}{Mixtures}: \omterm{def:g4:air-a-mixture-of-gases:air}{air}, sea water, milk.
\end{solution}

\begin{solution}{exo:g6:pure-substances-and-mixtures:2}
Homogeneous: salty water, tap water. Heterogeneous: granite, oil on
water, muesli.
\end{solution}

\begin{solution}{exo:g6:pure-substances-and-mixtures:3}
One particular kind of substance, with its own name and properties: for
instance water, salt, ethanol (also sugar, oxygen, iron).
\end{solution}

\begin{solution}{exo:g6:pure-substances-and-mixtures:4}
A fixed boiling temperature suggests a \omterm{def:g6:pure-substances-and-mixtures:pure-substance}{pure substance}; \qty{78}{\celsius}
is the boiling temperature of ethanol.
\end{solution}

\begin{solution}{exo:g6:pure-substances-and-mixtures:5}
It contains many \omterm{def:g6:pure-substances-and-mixtures:species}{chemical species} (water, sugars, acids, vitamins,
flavours): it is a \omterm{def:g6:pure-substances-and-mixtures:mixture}{mixture}. ``Pure'' on the carton only means that
nothing was added.
\end{solution}

\begin{solution}{exo:g6:pure-substances-and-mixtures:6}
No: a \omterm{def:g6:pure-substances-and-mixtures:pure-substance}{pure substance} boils at a fixed temperature. It is a \omterm{def:g6:pure-substances-and-mixtures:mixture}{mixture},
probably water with something dissolved in it, such as salty water.
\end{solution}

\begin{solution}{exo:g6:pure-substances-and-mixtures:7}
A filter could separate the muddy water (the grains stay on the paper)
and the juice with pulp (the pulp stays). It cannot separate salty water
(the salt is dissolved) nor oil from water (both liquids pass through;
they are separated by pouring off the top layer).
\end{solution}

\begin{solution}{exo:g6:pure-substances-and-mixtures:8}
$78.7 \div 10 = \qty{7.87}{g}$ for \qty{1}{cm^3}: iron.
\end{solution}

\begin{solution}{exo:g6:pure-substances-and-mixtures:9}
$\frac{357}{478} \approx 0.747$: about \qty{75}{\percent}.
\end{solution}

\begin{solution}{exo:g6:pure-substances-and-mixtures:10}
$2 \times 35 = \qty{70}{g}$ of salts. Percentage: $\frac{35}{1000} =
\qty{3.5}{\percent}$.
\end{solution}

\begin{solution}{exo:g6:pure-substances-and-mixtures:11}
Measure their boiling temperatures (each must stay fixed while boiling,
and the two must be equal), and weigh the same volume of each (equal
masses). Two equal constants, both fixed, point to the same \omterm{def:g6:pure-substances-and-mixtures:pure-substance}{pure
substance}.
\end{solution}

\begin{solution}{exo:g6:pure-substances-and-mixtures:12}
$\frac{10}{90 + 10} = \qty{10}{\percent}$ ethanol. It is a \omterm{def:g6:pure-substances-and-mixtures:mixture}{mixture}: it
has no fixed boiling temperature and appears nowhere in the table, which
lists only \omterm{def:g6:pure-substances-and-mixtures:pure-substance}{pure substances}.
\end{solution}

\begin{solution}{pb:g6:pure-substances-and-mixtures:1}
\textbf{1.} No. A \omterm{def:g6:pure-substances-and-mixtures:pure-substance}{pure substance} and a \omterm{def:g6:pure-substances-and-mixtures:homogeneous}{homogeneous mixture} can look
exactly alike: something must be measured.

\textbf{2.} Homogeneous: it is clear and its parts cannot be seen.

\textbf{3.} A constant: a boiling temperature that stays fixed while the
liquid boils (or the mass of a known volume, compared with a table).

\textbf{4.} A and B, whose boiling temperatures stay fixed, are \omterm{def:g6:pure-substances-and-mixtures:pure-substance}{pure
substances}. C is a \omterm{def:g6:pure-substances-and-mixtures:mixture}{mixture}.

\textbf{5.} A boils at \qty{100}{\celsius}: water. B boils at
\qty{78}{\celsius}: ethanol.

\textbf{6.} A: $49.8 \div 50.0 = \qty{0.996}{g}$. B: $39.5 \div 50.0 =
\qty{0.790}{g}$. C: $51.2 \div 50.0 = \qty{1.024}{g}$.

\textbf{7.} Yes: about \qty{1.0}{g} for water, \qty{0.79}{g} for
ethanol.

\textbf{8.} As the water boils away, the salt stays behind, so the
liquid that remains holds more and more salt; it then boils at a higher
temperature.

\textbf{9.} Liquid C holds a dissolved solid: it is a solution, a
\omterm{def:g6:pure-substances-and-mixtures:mixture}{mixture}, as its boiling already showed.

\textbf{10.} $100 \times 1.024 = \qty{102.4}{g}$; $\frac{3.5}{102.4}
\approx 0.034$: about \qty{3.4}{\percent} of the mass is salt.

\textbf{11.} $\frac{35}{1000} = \qty{3.5}{\percent}$: very close to
question 10.

\textbf{12.} \qty{3.5}{g} in \qty{100}{mL}, so $3.5 \times 10 =
\qty{35}{g}$ of salt per litre. With its saltiness and its boiling
behaviour, C is the sea water; A is distilled water and B ethanol.
\end{solution}
